\section{Introduction}
\label{sec:introduction}

Multi-turn conversation is now the default way people use language models, and a growing body of work shows that the same instructions produce worse output when they arrive gradually across a conversation instead of all at once in a single prompt \citep{laban2025llmslostmultiturnconversation,he2024multiifbenchmarkingllmsmultiturn,jia2026battleanotherprobingllms,li2025structflowbenchstructuredflowbenchmark}. Models forget constraints they had already satisfied, anchor on early assumptions, and struggle specifically when a later turn has to revise earlier content rather than simply add to it. This ``getting lost in conversation'' effect is large and consistent, but every study that documents it shares one property: the underlying task has an objectively verifiable answer, checkable by script rather than judgment. Word counts, exact phrases, code that passes a test, and dialogue turns that either recall or forget a stated fact all have a ground truth independent of the model that produced them.

Creative writing has no such ground truth. There is no single correct story for a given premise, constraints are frequently placed in tension with one another on purpose, and the qualities that make a piece of writing good, such as structure, originality, and how well it delivers on a genre, cannot be checked by script. It is also one of the most common ways people actually use language models in multi-turn conversation: a story, a screenplay outline, or a piece of collaborative fiction is rarely specified perfectly in one shot, and instead accumulates requirements, revisions, and new constraints turn by turn. Whether the get-lost effect documented on verifiable tasks shows up here, and if so, whether it looks the same, is unknown. A model that stops satisfying explicit requirements under incremental delivery is one failure mode; a model that keeps satisfying every requirement but produces a structurally worse or blander story is a different one, and the two carry different implications for how multi-turn creative-writing systems should be built and evaluated. Answering this also raises a problem the verifiable-task literature does not have to face: if the writing-quality measurements themselves come from an LLM judge, and that judge's notion of ``better writing'' does not track human preference, then any conclusion about creative degradation is only as trustworthy as the judge that produced it.

We study three research questions. \textbf{RQ1} (\cref{sec:rq1}): does delivering a creative-writing specification one piece at a time, rather than all at once, reduce how many of its explicit constraints the final story satisfies? \textbf{RQ2} (\cref{sec:rq2}): how does incremental delivery change writing quality across craft, structural coherence, originality, genre effectiveness, and characterization? \textbf{RQ3} (\cref{sec:rq3}): do any of those writing-quality differences survive when the comparison is restricted to outputs that satisfied the same fraction of constraints under both delivery conditions, which tests whether writing-quality differences persist when observed constraint adherence is equal?

To answer these questions we build a benchmark of 160 creative-writing tasks spanning six genres, each with a fully-specified instruction and a matched sharded version that distributes the task specification across five to nine conversational turns (\cref{sec:sharding}). We generate from six open-weight language models under both conditions, yielding 1{,}920 final outputs and 960 matched FULL/Sharded pairs (\cref{sec:models}), and score every output with a blinded automated judge against its task's atomic constraints and five creative-quality dimensions (\cref{sec:constraints}). Because that judge's creative-quality scores are the basis for RQ2 and RQ3, we separately compare a related pairwise LLM evaluator built on the same rubric concepts with human judgments, and audit both pairwise evaluators for stability under response-order reversal, reporting where automated judgment agrees with human preference and where it does not (\cref{sec:human_eval}).

\paragraph{Contributions.} This paper makes four contributions.
\begin{enumerate}
    \item \textbf{A paired FULL/Sharded creative-writing benchmark.} 160 tasks across six genres, each with a fully-specified instruction and a matched multi-turn sharded delivery designed to preserve the task specification, together with 1{,}278 extracted atomic constraints for automated scoring (\cref{sec:sharding,sec:constraints}).
    \item \textbf{Evidence that multi-turn degradation in creative writing is not one effect but at least two.} FULL delivery produces higher constraint adherence and better structure/coherence, but the structure/coherence gap persists even among pairs matched on observed constraint adherence, while originality shifts slightly toward Sharded (\cref{sec:experiments}) --- constraint forgetting alone does not explain the full gap.
    \item \textbf{Pairwise evidence on the limits of automated creative-writing judgment.} A 30-case human comparison of a related pairwise LLM evaluator, together with two pairwise order-swap audits, shows that automated judgments track human preference considerably better for constraint following than for holistic creative quality, and quantifies that gap directly rather than assuming the judge is reliable (\cref{sec:disc_human,sec:disc_robustness}).
    \item \textbf{A fully released benchmark, generation, and evaluation pipeline.} The task set, all 6{,}936 raw generations, every judge score, and the human-validation data are published for reuse, with full generation and scoring provenance preserved down to the individual seed and context-length event.
\end{enumerate}
